% ===========
%  DẪN NHẬP
% ===========
\OpenGrid

\ParallelChapterStar
{Preface}
{Dẫn nhập}
{``\emph{Preface}'' của sách không mang tính xã giao hay dẫn chuyện, mà giữ vai trò định hướng học thuật: xác lập mục tiêu, lập trường sư phạm và khung tiếp cận mà toàn bộ chương sau dựa vào. Cách dịch \emph{dẫn nhập} phản ánh đúng chức năng và phù hợp hơn với tinh thần lí luận - phân tích của sách.}

\ParallelSection
{Background}
{Bối cảnh}
{\emph{Bối cảnh}'' lột tả được nội dung phần này: trình bày hoàn cảnh và tiến trình chính sách - tổ chức dẫn tới sự hình thành ấn phẩm, không phải nền tảng lí luận hay nội dung học thuật của sách.}

\TriPara
{
    \EN{The National Council of Teachers of Mathematics (NCTM) has a long tradition of providing leadership and vision to support teachers in ensuring equitable mathematics learning of the highest quality for all students. In the three decades since the 1980 publication of \emph{An Agenda for Action}, NCTM has consistently advocated a coherent prekindergarten through grade 12 mathematics curriculum focused on mathematical problem solving. NCTM refined this message in \emph{Curriculum and Evaluation Standards for School Mathematics} (1989), which argued for a common core of mathematics for all students, with attention to the processes of problem solving, reasoning, connections, and communication. Continuing to provide leadership in 2000, NCTM issued \emph{Principles and Standards for School Mathematics} (\emph{Principles and Standards}) (2000a), which updated and elaborated on the 1989 recommendations in a set of five Content Standards and five Process Standards to include in every school mathematics program. In 2006, the Council's \emph{Curriculum Focal Points for Prekindergarten through Grade 8 Mathematics} (\emph{Curriculum Focal Points}) (NCTM 2006a) offered guidance on how to focus the mathematical content for each grade level from prekindergarten through grade 8 while reminding educators of the importance of consistently embedding the relevant mathematical processes throughout the content in every mathematical learning experience.}
}
{
    \VI{Hội đồng Quốc gia Giáo viên Toán học Hoa Kì (NCTM) có truyền thống lâu dài trong việc đưa ra định hướng và tầm nhìn nhằm hỗ trợ giáo viên bảo đảm mọi học sinh đều được tiếp cận nền giáo dục toán học công bằng, chất lượng cao nhất. Trong ba thập niên kể từ khi \emph{An Agenda for Action} được công bố năm 1980, NCTM luôn nhất quán ủng hộ chương trình toán học liên thông từ tiền mẫu giáo đến lớp 12, lấy giải quyết vấn đề toán học làm trọng tâm. NCTM làm rõ thêm thông điệp này trong \emph{Curriculum and Evaluation Standards for School Mathematics} (1989), trong đó bàn về cốt lõi toán học chung dành cho mọi học sinh, đồng thời chú ý đến quá trình giải quyết vấn đề, luận, kết nối và giao tiếp. Tiếp tục giữ vai trò dẫn dắt, năm 2000 NCTM ban hành \emph{Principles and Standards for School Mathematics} (\emph{Principles and Standards}) (2000a), cập nhật và triển khai chi tiết hơn khuyến nghị năm 1989 thành hệ thống gồm năm Chuẩn nội dung và năm Chuẩn quá trình cần có trong mọi chương trình toán ở nhà trường. Năm 2006, tài liệu \emph{Curriculum Focal Points for Prekindergarten through Grade 8 Mathematics} (\emph{Curriculum Focal Points}) (NCTM 2006a) của Hội đồng đưa ra định hướng về cách tập trung nội dung toán học ở từng cấp lớp từ tiền mẫu giáo đến lớp 8, đồng thời nhắc nhà giáo dục về tầm quan trọng của việc lồng ghép nhất quán các quá trình toán học phù hợp vào toàn bộ nội dung trong mọi trải nghiệm học toán.}
}
{
    \NOTE{FHSM được đặt vào tiến trình phát triển tư tưởng liên tục của NCTM: từ việc nhấn mạnh giải quyết vấn đề như trọng tâm của chương trình học toán (1980), đến việc xác lập phần cốt lõi toán học chung cho mọi học sinh cùng với sự chú ý rõ ràng đến các quá trình như giải quyết vấn đề, luận, kết nối và giao tiếp (1989), rồi hệ thống hóa định hướng ấy thành Chuẩn nội dung và Chuẩn quá trình cho toàn bộ chương trình toán ở nhà trường (2000), và tiếp đó là xác định trọng điểm nội dung theo từng cấp lớp đồng thời nhấn mạnh việc luôn lồng ghép quá trình toán học vào mọi trải nghiệm học toán (2006).

    Trong mạch tư tưởng đó, FHSM không xuất hiện như sáng kiến rời rạc, mà như bước tiếp tục làm nổi bật và đào sâu hướng đã được NCTM theo đuổi từ trước: việc học toán không thể chỉ được tổ chức quanh nội dung, mà còn phải được tổ chức quanh quá trình toán học cốt lõi. Điểm riêng của FHSM là đưa luận và hiểu vào vị trí trung tâm để từ đó đọc lại việc học toán trung học phổ thông.

    Cách định vị này giúp người đọc thấy rằng cuốn sách không nhằm bổ sung mảng nội dung mới cho chương trình, mà nhằm làm rõ điều gì phải giữ vai trò trung tâm trong trải nghiệm học toán, từ đó chi phối việc thiết kế chương trình học, dạy học và đánh giá.}
}

\TriPara
{
    \EN{Realizing a parallel need for focus and coherence at the high school level, although perhaps of a different kind from that addressed in \emph{Curriculum Focal Points} (NCTM 2006a), a 2006 NCTM task force recommended the development of a framework based on \emph{Principles and Standards} (NCTM 2000a) that would guide future work regarding high school mathematics. As a result, a writing group was appointed to produce the framework outlined in this publication. Concurrently, a planning committee was appointed to oversee the larger high school curriculum project, which included guiding the review process for this publication; planning its rollout, including pamphlets summarizing the vision presented in this publication for various audiences; and proposing a series of topic books that elaborate the publication's core messages.}
}
{
    \VI{Nhận thấy ở bậc trung học phổ thông cũng có nhu cầu tương ứng về sự tập trung và tính nhất quán, dù có thể khác về tính chất so với nhu cầu được giải quyết trong \emph{Curriculum Focal Points} (NCTM 2006a), năm 2006 một nhóm công tác của NCTM đã khuyến nghị xây dựng khuôn khổ dựa trên \emph{Principles and Standards} (NCTM 2000a) để định hướng công việc tiếp theo liên quan đến toán trung học phổ thông. Vì vậy, nhóm biên soạn được thành lập để soạn thảo khuôn khổ được phác thảo trong ấn phẩm này. Đồng thời, ban lập kế hoạch cũng được thành lập để giám sát dự án rộng hơn về chương trình toán trung học phổ thông. Dự án này bao gồm việc định hướng quá trình thẩm định ấn phẩm; lập kế hoạch phổ biến ấn phẩm, gồm cả các tập sách mỏng tóm tắt tầm nhìn được trình bày trong ấn phẩm này cho nhiều nhóm độc giả; và đề xuất loạt sách chuyên đề triển khai chi tiết hơn các thông điệp cốt lõi của ấn phẩm.}
}{\NOTE{}}

\ParallelSection
{Purpose}
{Mục đích}
{``\emph{Purpose}'' được dịch là ``\emph{Mục đích}'' vì phần này nhằm trình bày vai trò và ý đồ định hướng của FHSM, không phải các mục tiêu hay chỉ tiêu cần đạt.}

\TriPara
{
    \EN{This publication advocates that all high school mathematics programs focus on reasoning and sense making. The audience for this publication is intended to be everyone involved in decisions regarding high school mathematics programs, including formal decision makers within the system; people charged with implementing those decisions; and other stakeholders affected by, and involved in, those decisions.}
}
{
    \VI{Ấn phẩm này chủ trương rằng mọi chương trình toán trung học phổ thông cần lấy luận và hiểu làm trọng tâm. Đối tượng mà ấn phẩm hướng tới là tất cả những người tham gia vào các quyết định liên quan đến chương trình toán trung học phổ thông, bao gồm những người ra quyết định chính thức trong hệ thống; những người được giao trách nhiệm thực thi các quyết định đó; và những bên liên quan khác vừa chịu tác động bởi, vừa tham gia vào, các quyết định ấy.}
}
{
\NOTE{}
}

\TriPara
{
    \EN{A number of publications produced over the past few years have provided detailed analyses of the topics that should be addressed in each course of high school mathematics (cf. American Diploma Project 2004; College Board 2006, 2007; ACT 2007; Achieve 2007a, 2007b.) This publication takes a somewhat different approach, proposing curricular emphases and instructional approaches that make reasoning and sense making foundational to the content that is taught and learned. Along with the more-detailed content recommendations outlined in \emph{Principles and Standards}, this publication supplies a critical filter in examining any curriculum arrangement to ensure that the ultimate goals of the high school mathematics program are achieved.}
}
{
    \VI{Trong vài năm gần đây, một số ấn phẩm đã đưa ra những phân tích chi tiết về các chủ đề cần được đề cập trong từng khóa học toán trung học phổ thông (xem American Diploma Project 2004; College Board 2006, 2007; ACT 2007; Achieve 2007a, 2007b). Tuy nhiên, ấn phẩm này chọn cách tiếp cận hơi khác: đề xuất các trọng tâm chương trình và cách tiếp cận dạy học nhằm làm cho luận và hiểu trở thành nền tảng của nội dung được dạy và học. Cùng với những khuyến nghị nội dung chi tiết hơn được trình bày trong \emph{Principles and Standards}, ấn phẩm này cung cấp bộ lọc then chốt để xem xét bất kì cách tổ chức chương trình nào, nhằm bảo đảm đạt được các mục tiêu cuối cùng của chương trình toán trung học phổ thông.}
}
{
    \NOTE{Phần ``Mục đích'' xác lập rõ vai trò của FHSM: đây không phải là danh mục nội dung hay bộ chuẩn đầu ra mới, mà là khung định hướng nhằm đưa luận và hiểu trở thành nền tảng của nội dung được dạy và học ở bậc trung học phổ thông. Trong khi nhiều tài liệu khác tập trung trả lời câu hỏi ``cần dạy những chủ đề nào'', FHSM tập trung vào câu hỏi ``cần tổ chức việc dạy và học như thế nào để học sinh thực sự phát triển tư duy toán học''. Cụm \emph{bộ lọc then chốt} được dùng để nhấn mạnh chức năng của ấn phẩm như tiêu chí định hướng và đánh giá chương trình: với cùng một nội dung, cách tổ chức nào giúp đạt được năng lực tư duy toán học cốt lõi.}
}

\ParallelSection
{Organization of the Publication}
{Cấu trúc ấn phẩm}
{}

\TriPara
{
    \EN{The first section of the publication presents an overview of reasoning and sense making. Chapter 1 describes what constitutes reasoning and sense making in the mathematics classroom; why they should be considered as foundational for high school mathematics; and how they link with other mathematical processes, such as representation, communication, connections, and problem solving. Chapter 2 describes in more detail the mathematical reasoning habits that students should continue to acquire throughout their high school mathematics experiences, a general trajectory for how they develop, suggestions for how to promote them in the classroom, and an explanation of how reasoning and sense making fit into the larger picture of mathematical activity.}
}
{
    \VI{Phần thứ nhất của ấn phẩm trình bày tổng quan về luận và hiểu. Chương 1 mô tả điều gì cấu thành luận và hiểu trong lớp học toán; vì sao chúng cần được xem là nền tảng của toán trung học phổ thông; và cách chúng liên kết với các quá trình toán học khác như biểu diễn, giao tiếp, kết nối và giải quyết vấn đề. Chương 2 mô tả chi tiết hơn những thói quen luận toán học mà học sinh cần tiếp tục hình thành trong suốt trải nghiệm học toán trung học phổ thông; lộ trình tổng quát cho sự phát triển của các thói quen ấy; gợi ý về cách thúc đẩy chúng trong lớp học; cùng với giải thích về cách luận và hiểu nằm trong bức tranh rộng hơn của hoạt động toán học.}
}
{
\NOTE{}
}

\TriPara
{
    \EN{The second section of the publication demonstrates with examples how reasoning and sense making can be incorporated into the high school curriculum. Chapter 3 provides an overview, and chapters 4--8 describe how reasoning and sense making fit into five overarching areas of high school mathematics---number and measurement, algebraic symbols, functions, geometry, and statistics and probability.}
}
{
    \VI{Phần thứ hai của ấn phẩm minh họa bằng các ví dụ cách đưa luận và hiểu vào chương trình toán trung học phổ thông. Chương 3 trình bày tổng quan, còn các chương 4--8 mô tả cách luận và hiểu gắn với năm lĩnh vực bao quát của toán trung học phổ thông---số và đo lường, kí hiệu đại số, hàm số, hình học, thống kê và xác suất.}
}
{
\NOTE{}
}

\TriPara
{
    \EN{The final section discusses issues in implementing reasoning and sense making across the high school mathematics program. Chapter 9 focuses on how to provide equitable opportunities for all students to engage in reasoning and sense making. Chapter 10 addresses the importance of coherent expectations regarding curriculum, instruction, and assessment in promoting reasoning and sense making. Finally, chapter 11 presents questions to consider as stakeholders work together to improve high school mathematics education.}
}
{
    \VI{Phần cuối thảo luận những vấn đề liên quan đến việc triển khai luận và hiểu xuyên suốt chương trình toán trung học phổ thông. Chương 9 tập trung vào cách tạo cơ hội công bằng để mọi học sinh tham gia vào luận và hiểu. Chương 10 bàn về tầm quan trọng của những kì vọng nhất quán đối với chương trình học, dạy học và đánh giá trong việc thúc đẩy luận và hiểu. Cuối cùng, chương 11 trình bày những câu hỏi cần cân nhắc khi các bên liên quan cùng hợp tác cải thiện giáo dục toán trung học phổ thông.}
}
{
    \NOTE{Ba phần của ấn phẩm tạo thành một lộ trình có chủ ý: từ việc làm rõ luận và hiểu là gì, vì sao chúng giữ vai trò nền tảng; sang minh họa cách chúng vận hành trong các lĩnh vực nội dung chính của toán trung học phổ thông; và cuối cùng là thảo luận những vấn đề triển khai ở cấp chương trình và hệ thống giáo dục. Cách tổ chức này cho thấy FHSM không chỉ dừng ở lí luận hay ví dụ lớp học, mà còn hướng tới việc tác động đồng bộ đến thiết kế chương trình, dạy học và đánh giá.}
}

\TriPara
{
    \EN{For the convenience of the reader, an annotated bibliography is included that describes pertinent research works that underlie this publication, organized by chapter.}
}
{
    \VI{Để thuận tiện cho người đọc, ấn phẩm có kèm phần Thư mục chú giải, mô tả những công trình nghiên cứu liên quan làm nền tảng cho ấn phẩm này và được sắp xếp theo từng chương.}

}
{
    \NOTE{}
}

\CloseGrid